\documentclass{article}
\usepackage{graphicx} % Required for inserting images
\usepackage{amsmath}
\title{Multifit-TB}
\author{Ruben }
\date{May 2026}

\begin{document}

\maketitle



The measured discrete waveform $w(t)$ is modeled as a scaled and time-shifted version of a reference pulse template $P(t)$:
\begin{equation}
w(t) = A \cdot P(t - \Delta t) + \epsilon(t),
\end{equation}
where:
\begin{itemize}
    \item $A$ is the pulse amplitude,
    \item $\Delta t$ is the time shift,
    \item $\epsilon(t)$ represents measurement noise.
\end{itemize}

The goal is to estimate $(A, \Delta t)$ for each waveform.


The model is nonlinear due to the time shift $\Delta t$. To obtain a tractable solution, a first-order Taylor expansion of the pulse template is performed around $\Delta t = 0$:
\begin{equation}
P(t - \Delta t) \approx P(t) - \Delta t \cdot P'(t),
\end{equation}
where $P'(t)$ denotes the derivative of the pulse shape.

Substituting into the model:
\begin{equation}
w(t) \approx A \cdot P(t) - A \cdot \Delta t \cdot P'(t).
\end{equation}

Defining the auxiliary parameter
\begin{equation}
C = -A \cdot \Delta t,
\end{equation}
the model becomes linear:
\begin{equation}
w(t) \approx A \cdot P(t) + C \cdot P'(t).
\end{equation}


The optimal parameters $(A, C)$ are obtained by minimizing the quadratic cost function:
\begin{equation}
\chi^2(A, C) = \sum_{i=1}^{N} \left[ w(t_i) - A P(t_i) - C P'(t_i) \right]^2.
\end{equation}

Setting the derivatives of $\chi^2$ with respect to $A$ and $C$ to zero yields the normal equations:
\begin{align}
\sum_i P(t_i)\, w(t_i) &= A \sum_i P(t_i)^2 + C \sum_i P(t_i) P'(t_i), \\
\sum_i P'(t_i)\, w(t_i) &= A \sum_i P(t_i) P'(t_i) + C \sum_i P'(t_i)^2.
\end{align}


Introducing the following scalar products:
\begin{align}
Ap   &= \sum_i w(t_i) P(t_i), \\
Ad   &= \sum_i w(t_i) P'(t_i), \\
PP   &= \sum_i P(t_i)^2, \\
PdP  &= \sum_i P(t_i) P'(t_i), \\
dPdP &= \sum_i P'(t_i)^2,
\end{align}
the system can be written in matrix form:
\begin{equation}
\begin{bmatrix}
PP & PdP \\
PdP & dPdP
\end{bmatrix}
\begin{bmatrix}
A \\
C
\end{bmatrix}
=
\begin{bmatrix}
Ap \\
Ad
\end{bmatrix}.
\end{equation}


The solution is obtained by inversion of the $2 \times 2$ system:
\begin{equation}
\Delta = PP \cdot dPdP - (PdP)^2,
\end{equation}
\begin{align}
A &= \frac{Ap \cdot dPdP - Ad \cdot PdP}{\Delta}, \\
C &= \frac{Ad \cdot PP - Ap \cdot PdP}{\Delta}.
\end{align}


Using the relation $C = -A \Delta t$, the time shift is given by:
\begin{equation}
\Delta t = -\frac{C}{A}.
\end{equation}



The linearization is only valid for small $\Delta t$. Therefore, the procedure is applied iteratively:

\begin{enumerate}
    \item Initialize $\Delta t^{(0)} = 0$.
    \item At iteration $k$, evaluate the shifted template:
    \begin{equation}
    P^{(k)}(t) = P(t - \Delta t^{(k)}), \quad
    P'^{(k)}(t) = P'(t - \Delta t^{(k)}).
    \end{equation}
    \item Solve the linear system for $(A^{(k)}, C^{(k)})$.
    \item Update the time shift:
    \begin{equation}
    \Delta t^{(k+1)} = \Delta t^{(k)} - \frac{C^{(k)}}{A^{(k)}}.
    \end{equation}
    \item Repeat until convergence or for a fixed number of iterations.
\end{enumerate}

This procedure corresponds to a Gauss--Newton minimization of the nonlinear least-squares problem.


The method can be interpreted as projecting the waveform onto two basis functions:
\begin{itemize}
    \item $P(t)$, which captures the amplitude,
    \item $P'(t)$, which captures sensitivity to time shifts.
\end{itemize}

The coefficient $A$ determines the signal strength, while the ratio $C/A$ encodes the temporal misalignment between the waveform and the template.



\end{document}
